\documentclass[10pt,a4paper]{amsart}
\usepackage[margin=19mm]{geometry}
\usepackage{amsmath,amssymb,mathtools}
\usepackage[scr=rsfs,cal=euler]{mathalfa}
\usepackage{xfrac}
\usepackage[unicode,hidelinks,bookmarks=false]{hyperref}
\newcommand{\ie}{i.e.\ }
\newcommand{\cf}{cf.\ }
\newcommand{\resp}{respectively\ }
\newcommand{\Subsection}{\subsection}
\providecommand{\nobreakdash}{\nobreak\hbox{-}}
\setlength{\emergencystretch}{2em}
\multlinegap0pt
\usepackage[shortlabels]{enumitem}
\usepackage{amssymb}
\usepackage{mathtools}
\usepackage{xcolor}
\newtheorem{theorem}{Theorem}
\newtheorem{proposition}{Proposition}
\newtheorem{lemma}{Lemma}
\newtheorem{corollary}{Corollary}
\usepackage{cleveref}
\crefname{theorem}{Theorem}{Theorems}
\crefname{proposition}{Proposition}{Propositions}
\crefname{lemma}{Lemma}{Lemmas}
\crefname{corollary}{Corollary}{Corollarys}
\newcommand{\Ext}{\operatorname{Ext}}
\newcommand{\Hom}{\operatorname{Hom}}
\newcommand{\modu}{\operatorname{mod}}
\newcommand{\perpA}{{}^{\perp}\!A}
\newcommand{\rad}{\operatorname{rad}}
\newcommand{\stableHom}{\underline{\Hom}}
\title{The Auslander-Reiten conjecture for algebras with radical cube zero}
\author{Xiaojin Zhang, Panyue Zhou}
\date{Source: arXiv:2609.08679v1 (CC BY 4.0)}
\begin{document}
\maketitle

\noindent\emph{Source and licence.} The Auslander-Reiten conjecture for algebras with radical cube zero. Version arXiv:2609.08679v1 (CC BY 4.0), distributed by arXiv under the Creative Commons Attribution 4.0 International licence, http://creativecommons.org/licenses/by/4.0/, as recorded in the arXiv licence metadata for this exact version and shown on the abstract page https://arxiv.org/abs/2609.08679v1. Full licence notice: \texttt{licenses/arxiv-2609.08679-CC-BY-4.0.txt}.\par\smallskip

\noindent\textit{Editorial scope.} Counted unit: the article's theorem thm:main (source0.tex lines 142--152). The article's complete original proof of it is delivered with it (source0.tex lines 554--599, one \texttt{proof} environment of 46 lines, which the article places away from the statement). No mathematics is written, repaired or re-derived: the statement and the proof are the article's own lines, in the article's own notation, and no retained line is substituted or re-flowed.  Retained uncounted as the source-local context the proof needs: lines 253--267; lines 303--313; lines 357--366; lines 418--420; lines 507--512; lines 521--530.  Nothing was omitted from any retained span, so the package carries no omission record.  No retained line contains a citation, so the wrapper carries no bibliography and no bibliography entry.  No retained line contains a figure environment, an image-inclusion command or drawing code, so no image is delivered and the package declares no asset dependency.  Editorial formatting only, in addition: an AMS \texttt{amsart} preamble that restates the article's own declarations and macros used by the retained lines, namely the environments \texttt{theorem}, \texttt{proposition}, \texttt{lemma}, \texttt{corollary} on separate counters, together with the standard packages those lines need.  The retained environments print in this document as Theorem 1, Proposition 1, Lemma 1, Corollary 1 in order of appearance, whereas the article prints the same items with the numbering of its own full document.  The complete native source of the article as distributed in the arXiv e-print for this exact version is bundled unchanged as \texttt{sources/209734/source0.tex}.
\begin{proposition}\label{prop:chain}
Let $L\ge1$ and let $M\in\perpA$ be non-projective. Assume that
$\Ext_A^q(M,M)=0$ for $1\le q\le L$. There exist indecomposable
non-projective modules $X_0,\ldots,X_{L-1}$ such that:
\begin{enumerate}[label=\textnormal{(\roman*)},leftmargin=2.2em]
 \item {$X_0\mid\Omega_A M$ and $X_{i+1}\mid\Omega_A X_i$.}
 \item {$X_i\in\perpA$ and $\Ext_A^q(X_i,X_i)=0$ for $1\le q\le L$.}
 \item if $j>i$, then
 \[
 \stableHom_A(X_j,X_i)=0=\Ext_A^1(X_j,X_i).
 \]
 \item the $X_i$ are pairwise non-isomorphic.
\end{enumerate}
If $J^3=0$, then $J^2X_i=0$ for every $i$.
\end{proposition}

\begin{lemma}\label{lem:quotient}
Let $I$ be a two-sided ideal of $A$, put $B=A/I$, and let $X,Y$ be
$B$-modules. Regard a $B$-module as an $A$-module through the quotient map
$A\twoheadrightarrow B$. This restriction of scalars, often called
inflation, induces an injection
\[
 \Ext_B^1(X,Y)\hookrightarrow\Ext_A^1(X,Y).
\]
Moreover, $\stableHom_A(X,Y)=0$ implies
$\stableHom_B(X,Y)=0$.
\end{lemma}

\begin{lemma}\label{lem:F}
The functor $F:B\text{-}\modu\to H\text{-}\modu$ is full, with
\[
 \ker\bigl(\Hom_B(X,Y)\longrightarrow\Hom_H(FX,FY)\bigr)
 \cong\Hom_S(X/\mathfrak rX,\mathfrak rY).
\]
It preserves indecomposability and projectivity and reflects
isomorphisms. For an indecomposable $X$, the module $FX$ is
projective if and only if $X$ is projective.
\end{lemma}

\begin{equation}\label{eq:ext-implication}
 \Ext_B^1(X,Y)=0\quad\Longrightarrow\quad\Ext_H^1(FX,FY)=0.
\end{equation}

\begin{corollary}\label{cor:hom-transfer}
Let $X$ be an indecomposable non-projective $B$-module. Then
\[
 \stableHom_B(X,Y)=0\quad\Longrightarrow\quad\Hom_H(FX,FY)=0.
\]
\end{corollary}

\begin{lemma}\label{lem:rank}
Let $H$ be a finite-dimensional hereditary algebra with $N$ simple
modules up to isomorphism. Suppose that non-zero modules
$Y_0,\ldots,Y_{r-1}$ satisfy
\[
 \Ext_H^1(Y_j,Y_i)=0~(j\ge i),
 \quad\Hom_H(Y_j,Y_i)=0~(j>i).
\]
Then $r\le N$.
\end{lemma}

\begin{theorem}\label{thm:main}
Let $A$ be a split finite-dimensional $k$-algebra, $J=\rad A$,
and let $s$ be the number of isomorphism classes of simple
$A$-modules. Assume that $J^3=0$. For $M\in\perpA$, the following
conditions are equivalent:
\begin{enumerate}[label=\textnormal{(\roman*)},leftmargin=2.2em]
 \item $M$ is projective.
 \vspace{2mm}
 \item $\Ext_A^q(M,M)=0$ for $1\le q\le 3s+1$.
\end{enumerate}
\end{theorem}

\begin{proof}[\bf \emph{Proof of \cref{thm:main}}]
Projectivity implies the required vanishing. For the converse, the
assertion is invariant under Morita equivalence. Indeed, an exact
equivalence preserves extension groups, projectives and the number
$s$ of simple modules. It sends the regular module to a projective
generator, whose additive closure is the full subcategory of
projectives, so it preserves the condition defining $\perpA$.
It also preserves radical filtrations. Hence $J^3=0$ is preserved.
We may therefore assume that $A$ is basic.

Suppose that $M\in\perpA$ is non-projective and
$\Ext_A^q(M,M)=0$ for $1\le q\le L$, where $L=3s+1$.
Take $X_0,\ldots,X_{L-1}$ from \cref{prop:chain}, and put
\[
 B=A/J^2.
\]
The modules $X_i$ are annihilated by $J^2$, so they are $B$-modules.
They remain indecomposable and pairwise non-isomorphic, because
the morphisms between $B$-modules are unchanged by inflation to $A$.
By \cref{lem:quotient},
\begin{equation}\label{eq:B-orthogonal}
 \Ext_B^1(X_j,X_i)=0\quad(j\ge i),
 \quad\stableHom_B(X_j,X_i)=0\quad(j>i).
\end{equation}

There are exactly $s$ isomorphism classes of indecomposable
projective $B$-modules. Since the $X_i$ are pairwise non-isomorphic,
at most $s$ of them are projective over $B$. Retain all remaining
terms, with indices $i_0<\cdots<i_{r-1}$, and set
\[
 Y_a=F(X_{i_a})\quad(0\le a<r).
\]
Then $r\ge L-s=2s+1$. By \cref{lem:F}, each $Y_a$ is
indecomposable and non-projective. Formulas
\eqref{eq:ext-implication} and \eqref{eq:B-orthogonal} give
\[
 \Ext_H^1(Y_b,Y_a)=0\quad(b\ge a),
\]
including when the original $X_{i_a}$ is simple. Similarly,
\cref{cor:hom-transfer} gives
\[
 \Hom_H(Y_b,Y_a)=0\quad(b>a).
\]
The separated algebra $H$ has $2s$ simple modules. Therefore
\cref{lem:rank} gives $r\le2s$, contradicting $r\ge2s+1$.
\end{proof}

\end{document}
